
\documentclass{article}
\usepackage{colortbl}
\usepackage{makecell}
\usepackage{multirow}
\usepackage{supertabular}

\begin{document}

\newcounter{utterance}

\centering \large Interaction Transcript for game `cladder', experiment `full\_v1.5\_default', episode 8830 with selfrefine{-}llama{-}3.1{-}8b{-}instruct.
\vspace{24pt}

{ \footnotesize  \setcounter{utterance}{1}
\setlength{\tabcolsep}{0pt}
\begin{supertabular}{c@{$\;$}|p{.15\linewidth}@{}p{.15\linewidth}p{.15\linewidth}p{.15\linewidth}p{.15\linewidth}p{.15\linewidth}}
    \# & \multicolumn{2}{c}{Player} && \multicolumn{2}{c}{Game Master} \\
    \hline

    \theutterance \stepcounter{utterance}  
    & & & \multicolumn{4}{p{0.6\linewidth}}{
        \cellcolor[rgb]{0.9,0.9,0.9}{
            \makecell[{{p{\linewidth}}}]{
                \texttt{\tiny{[P1$\langle$GM]}}
                \texttt{You are an expert in causal inference. The following question is not a typical commonsense query, but rather a meticulously designed question created by a professor specializing in causal inference, intended to assess the students' mastery of the course content.} \\
\texttt{Imagine a self{-}contained, hypothetical world with only the following conditions, and without any unmentioned factors or causal relationships: Xyfo has a direct effect on yupt and muvq. Zupj has a direct effect on yupt. Yupt has a direct effect on muvq. Xyfo is unobserved. Method 1: We look at how yupt correlates with muvq case by case according to xyfo. Method 2: We look directly at how yupt correlates with muvq in general. To understand how yupt affects muvq, is it more correct to use the Method 1 than Method 2?} \\
\texttt{Start your answer with 'Yes' or 'No', followed by additional reasoning or evidence to support your explanation.} \\
            }
        }
    }
     \\ \\

    \theutterance \stepcounter{utterance}  
    & \multicolumn{4}{p{0.6\linewidth}}{
        \cellcolor[rgb]{0.561,0.737,0.561}{
            \makecell[{{p{\linewidth}}}]{
                \texttt{\tiny{[P1$\rangle$GM]}}
                \texttt{No, it is not more correct to use Method 1 than Method 2 to understand how yupt affects muvq.} \\
\\ 
\texttt{In this scenario, we have a complex causal structure with multiple direct effects and a common effect (yupt). The presence of Xyfo, which is unobserved, creates a problem of confounding. When we look at the correlation between yupt and muvq case by case according to Xyfo (Method 1), we are essentially adjusting for the effect of Xyfo on both yupt and muvq. However, since Xyfo is unobserved, we cannot accurately adjust for its effect.} \\
\\ 
\texttt{In contrast, Method 2 looks directly at the correlation between yupt and muvq in general, without adjusting for any intermediate variables. However, this approach is still valid because yupt is a common effect of both Zupj and Xyfo, and its effect on muvq is direct. By looking at the general correlation between yupt and muvq, we are essentially estimating the total effect of yupt on muvq, which includes both the direct effect and the indirect effect mediated by Xyfo.} \\
\\ 
\texttt{In this case, the total effect of yupt on muvq is the sum of the direct effect and the indirect effect mediated by Xyfo. Since Xyfo is unobserved, we cannot estimate the indirect effect separately, but we can still estimate the total effect by looking at the general correlation between yupt and muvq.} \\
\\ 
\texttt{Therefore, Method 2 is a more correct approach than Method 1 to understand how yupt affects muvq, as it provides a more comprehensive estimate of the total effect of yupt on muvq.} \\
            }
        }
    }
    & & \\ \\

    \theutterance \stepcounter{utterance}  
    & & & \multicolumn{2}{p{0.3\linewidth}}{
        \cellcolor[rgb]{0.9,0.9,0.9}{
            \makecell[{{p{\linewidth}}}]{
                \texttt{\tiny{[GM$|$GM]}}
                \texttt{No} \\
            }
        }
    }
    & & \\ \\

    \theutterance \stepcounter{utterance}  
    & & & \multicolumn{2}{p{0.3\linewidth}}{
        \cellcolor[rgb]{0.9,0.9,0.9}{
            \makecell[{{p{\linewidth}}}]{
                \texttt{\tiny{[GM$|$GM]}}
                \texttt{yes} \\
            }
        }
    }
    & & \\ \\

    \theutterance \stepcounter{utterance}  
    & & & \multicolumn{2}{p{0.3\linewidth}}{
        \cellcolor[rgb]{0.9,0.9,0.9}{
            \makecell[{{p{\linewidth}}}]{
                \texttt{\tiny{[GM$|$GM]}}
                \texttt{game\_result = LOSE} \\
            }
        }
    }
    & & \\ \\

\end{supertabular}
}

\end{document}
